
\begin{table}[ht]
\centering
\footnotesize
\setlength{\tabcolsep}{4pt}
\begin{tabular}{lllcccc}
\toprule
\makecell{\textbf{Family}\\\textbf{(mod 256)}} &
\textbf{Example program} & 
\textbf{Its output} &
\makecell{\textbf{Smallest}\\\textbf{size}} &
\makecell{\textbf{Earliest}\\\textbf{round}} &
\makecell{\textbf{Size at}\\\textbf{earliest round}} &
\makecell{$\mathbb{E}[\text{first round}]$\\\textbf{(univ.\ prior)}} \\
\midrule
Arithmetic & \texttt{S+[.++]}            & $1,3,5,7,9,\ldots$   & 99k  & 0 & 6.2M & $\approx 105$ \\
Fibonacci  & \texttt{S,[[.C>.C>]}        & $1,1,2,3,5,\ldots$   & 3.1M & 512 & 6.2M & $>53{,}000$ \\
Geometric  & \texttt{S+[.L>]}            & $1,3,9,27,81,\ldots$ & 3.1M & 256 & 3.1M & $>53{,}000$ \\
Quadratic  & \texttt{S,.[<C>{}>VX<RX++]}   & $9,25,59,111,\ldots$ & 3.1M & 512 & 6.2M & $>53{,}000$ \\
Cubic      & \texttt{S+[[-.L>L>-]-]}     & $0,254,236,74,\ldots$ & 495k & 512 & 6.2M & $>53{,}000$ \\
\bottomrule
\end{tabular}
\vspace{6pt}
\caption{Program families with recognizable mathematical structure discovered by the generator during self-play.  
\emph{Smallest size} reports the smallest model that generated at least a single member of the family during training. 
\emph{Earliest round} and \emph{size at earliest round} report the earliest round such a member first appears during training, and the model size at which this first round occurs.
Several of the discovered programs use the added primitives (\cref{apx:turing-machine}).
These shorten the programs and so raise their probability under baseline uniform sampling too.
The baseline draws from the same augmented alphabet and still finds none of these
families, so the gap between the two columns is attributable to self-play. See \cref{apx:emergent-math-structure} for additional experimental details.}
\label{tab: discovered-math-structure}
\end{table}
